% =====================================================================
% Capítulo 8 — Conclusão
% Dissertação MAS-Hunt (ABNT, PT-BR)
% Síntese + contribuições REFORMULADAS (não reivindicar ganho de detecção
% da governança) + trabalhos futuros. Veredito honesto de aggregate-v2.md.
% Chaves \cite verificadas: zhan2024injecagent, NEURIPS2024_eb113910,
% huang2024resilience, ali2023huntgptintegratingmachinelearningbased.
% =====================================================================

\chapter{Conclusão}
\label{cap:conclusao}

\section{Síntese do Trabalho}
\label{sec:concl-sintese}

Esta dissertação apresentou o MAS-Hunt (\textit{Multi-Agent System for Threat
Hunting}), uma arquitetura de governança multiagente em três camadas --- Conselho
de Governança, Gerentes e Trabalhadores/Caçadores --- para caça proativa a
ameaças em ambientes Windows instrumentados com Sysmon e Elastic Security,
materializada como um \textit{plugin} do Claude Code que orquestra subagentes
nativos por meio de habilidades, comandos de barra e \textit{hooks}. O sistema
incorpora cinco mecanismos de governança --- integridade de memória (M1),
validação cruzada (M2), resistência a injeção via \textit{framework} estruturado
de análise (M3), monitoramento comportamental (M4) e quarentena (M5) --- e foi
avaliado em um experimento controlado e reprodutível que comparou quatro condições
(C1-full, C1-nohardening, C2-naive e C3-Elastic) sobre janelas de telemetria e
verdade de referência idênticas, sob injeção adversarial de logs em múltiplos
modos e níveis de sofisticação.

A evidência experimental foi relatada com integridade, inclusive em seu achado
contrário à expectativa inicial. Três conclusões quantitativas estruturam o
trabalho. Primeira: um analista baseado em modelo de linguagem supera as regras
estáticas de SIEM por aproximadamente $4{,}5$ vezes em F2, com intervalos de
confiança amplamente separados, em um ganho dirigido pela revocação. Segunda: a
camada de governança multiagente \emph{não} melhora o F2 de detecção sobre um
modelo de linguagem de passagem única --- C1-full ($0{,}675$) e C2-naive
($0{,}672$) são estatisticamente indistinguíveis. Terceira: o valor distintivo e
mensurável da governança é a resistência à injeção adversarial --- somente
C1-full detectou e sinalizou conteúdo adversarial injetado, em 407 janelas, contra
zero nas três demais condições.

\section{Resposta à Questão de Pesquisa}
\label{sec:concl-resposta}

A questão de pesquisa indagou como projetar uma arquitetura multiagente para caça
a ameaças que opere sobre telemetria corporativa viva e demonstre resiliência
contra ataques voltados aos próprios agentes. A resposta sustentada pela evidência
é que o ganho de \emph{detecção} advém da interposição de um analista de modelo de
linguagem entre a telemetria e a regra estática --- e não da camada de governança
em si --- ao passo que o ganho de \emph{segurança} advém especificamente dos
mecanismos de endurecimento (M1 e M3), que transformam a manipulação adversarial
em evidência observável e auditável. Em outras palavras, a arquitetura multiagente
governada não detecta melhor do que um único agente competente, mas é a única,
entre as condições avaliadas, capaz de perceber e declarar que foi atacada ---
propriedade ausente tanto no agente isolado quanto na linha de base SIEM.

\section{Contribuições (Reformuladas)}
\label{sec:concl-contribuicoes}

À luz dos resultados, as contribuições do trabalho são enunciadas em sua forma
honesta e reformulada, distinta da expectativa original de que a governança
elevaria a acurácia de detecção:

\begin{enumerate}
    \item \textbf{Evidência quantitativa da superioridade do analista agêntico
    sobre regras de SIEM.} Sob telemetria empresarial realista e contaminação
    adversarial, demonstrou-se um ganho de $\approx 4{,}5\times$ em F2 das
    condições baseadas em modelo de linguagem sobre as regras padrão do Elastic
    Security, dirigido por uma recuperação de revocação que reduz os falsos
    negativos de $\approx 88\%$ para faixas operacionalmente úteis. Esta é a
    contribuição de detecção, e ela é arquiteturalmente atribuível ao analista de
    modelo de linguagem, e não à governança \cite{ali2023huntgptintegratingmachinelearningbased}.
    \item \textbf{Uma arquitetura de governança que adiciona resistência à injeção
    e auditabilidade sem custo de detecção.} Os mecanismos M1--M5, e em especial
    M1 e M3, conferem a C1-full a capacidade --- exclusiva entre as condições
    avaliadas --- de detectar e sinalizar a contaminação adversarial de seu
    próprio contexto analítico (407 janelas sinalizadas contra 0), mantendo F2 de
    detecção estatisticamente equivalente ao do agente não governado. A
    contribuição de governança é, portanto, de segurança e de observabilidade, e
    não de acurácia \cite{zhan2024injecagent, NEURIPS2024_eb113910, huang2024resilience}.
    Uma reavaliação adversarial complementar, sob um ataque de fabricação de
    evidências mais forte e controlando um vazamento de rótulo da telemetria,
    confirmou esse valor de auditabilidade --- detecção de injeção em $100\%$ das
    janelas maliciosas e nenhum comprometimento silencioso --- de forma robusta
    entre níveis de modelo trabalhador, sem que daí decorra qualquer vantagem de
    acurácia de detecção.
    \item \textbf{Uma metodologia de avaliação bidimensional e um arcabouço
    reprodutível.} O trabalho entrega uma matriz de avaliação adversarial
    --- corpus de 185 \textit{playbooks} LOLBin, injeções em quatro modos e três
    níveis, gatilho de saúde de telemetria, extração priorizada de EID1 e
    pontuação determinística por janela com intervalos de confiança por
    \textit{bootstrap} --- que pontua os agentes não apenas pelo acerto da
    classificação, mas também pela capacidade de detectar que foram atacados. Essa
    dupla dimensão de avaliação é, ela própria, uma contribuição metodológica
    transferível para a pesquisa em agentes de segurança.
\end{enumerate}

\section{Trabalhos Futuros}
\label{sec:concl-futuros}

A evidência reunida e as limitações reconhecidas apontam cinco direções de
continuidade:

\begin{enumerate}
    \item \textbf{Saturar a telemetria discriminativa.} Reconfigurar a coleta de
    Sysmon e complementá-la com fontes adicionais (ETW, EDR de
    \textit{endpoint}, telemetria de PowerShell e de rede) para elevar o teto
    estrutural de revocação imposto pela captura esparsa de EID1, isolando assim
    o efeito da governança da escassez de sinal. Só com a telemetria saturada será
    possível afirmar com segurança se a igualdade de F2 entre C1-full e C2-naive
    persiste.
    \item \textbf{Estressar adversarialmente a detecção, e não apenas a
    observabilidade.} Conceber \textit{payloads} de injeção que efetivamente
    alterem a evidência consumida pelo classificador (e não apenas o instruam a
    ignorá-la), de modo a testar se a resistência da governança se traduz em
    \emph{melhor acurácia sob ataque forte} --- regime em que o agente de passagem
    única tenderia a falhar silenciosamente.
    \item \textbf{Reduzir a taxa de falsos positivos.} Incorporar calibração,
    agrupamento de alertas, validação determinística adicional e aprendizado com
    \textit{feedback} do analista, transpondo o ponto de operação otimizado para
    F2 a um ponto operacionalmente viável em SOC.
    \item \textbf{Ampliar o escopo de avaliação.} Estender o corpus a técnicas
    pós-exploração além de LOLBins e a ambientes Linux, macOS e nuvem pública, com
    validação longitudinal em redes de maior escala e diversidade de
    \textit{background}.
    \item \textbf{Separar o efeito do modelo do efeito da governança.} Reavaliar a
    arquitetura com diferentes modelos de linguagem subjacentes e medir custos
    operacionais (latência de triagem, consumo de \textit{tokens}, custo
    financeiro e tempo até a validação de evidência), de modo a caracterizar o
    \textit{trade-off} entre robustez de governança e custo computacional.
\end{enumerate}

\section{Considerações Finais}
\label{sec:concl-finais}

O MAS-Hunt não elimina a necessidade de SIEM, de telemetria confiável ou de
validação humana; ele os pressupõe e sobre eles opera. A principal lição deste
trabalho não é a que se esperava no início: a governança multiagente não tornou o
sistema um detector mais acurado. O que ela tornou o sistema foi um detector mais
\emph{honesto} sob ataque --- capaz de declarar, sob manipulação adversarial, que
seu próprio raciocínio foi alvo de contaminação, em vez de produzir uma
classificação aparentemente correta sem qualquer rastro do ataque que sofreu. Em
segurança cibernética, na qual o adversário explora justamente a confiança cega no
processo de decisão, essa diferença --- entre errar (ou acertar) em silêncio e
deixar evidência auditável da manipulação --- é precisamente o que separa um
agente conveniente de um agente defensável. É essa propriedade, e não um número de
acurácia superior, que a evidência experimental desta dissertação sustenta como a
contribuição de segurança do MAS-Hunt.
